% Part: normal-modal-logic
% Chapter: tableaux
% Section: simple-S5

\documentclass[../../../include/open-logic-section]{subfiles}

\begin{document}

\olfileid[pt]{nml}{tab}{s5}

\olsection{\usetoken{P}{tableau} simples para \Log{S5}}

\Log{S5} é correta e completa em relação à classe dos modelos universais,
O.s.\@, modelos nos quais todo mundo é acessível a partir de todo mundo.
Em modelos universais, a relação de acessibilidade não faz diferença:
``há um mundo~$w$ em que $\mSat{M}{!A}[w]$'' é verdadeira se, e somente
se, há um tal $w$ acessível a partir de~$u$. Assim, em \Log{S5}, podemos
definir modelos simplesmente como um conjunto de mundos e uma valoração
~$V$. Isso sugere que também devemos poder simplificar as regras de
!!{tableau}. No caso geral, tomamos como prefixos sequências de inteiros
positivos, de modo a registrar quais prefixos nomeiam mundos acessíveis a
partir de outros: $\sigma.n$ nomeia um mundo acessível a partir de
~$\sigma$. Mas, em \Log{S5}, qualquer mundo é acessível a partir de
qualquer outro; portanto, não é necessário manter esse registro. Em vez
disso, podemos usar inteiros positivos como prefixos. As regras
simplificadas são apresentadas em \olref{tab:rules-S5}.

\begin{table}
  \begin{center}
    \def\arraystretch{3}\def\fCenter{}
    \begin{tabular}{|c|c|}
    \hline
    \iftag{prvBox}{
      \AxiomC{\sFmla{\True}{\Box !A}[n]}
      \RightLabel{\TRule{\True}{\Box}}
      \UnaryInfC{\sFmla{\True}{!A}[m]}
      \DisplayProof
      &
      \AxiomC{\sFmla{\False}{\Box !A}[n]}
      \RightLabel{\TRule{\False}{\Box}}
      \UnaryInfC{\sFmla{\False}{!A}[m]}
      \DisplayProof\\
      $m$ é usado & $m$ é novo
      \\[1ex]
      \hline}{}
    \iftag{prvDiamond}{
      \AxiomC{\sFmla{\True}{\Diamond !A}[n]}
      \RightLabel{\TRule{\True}{\Diamond}}
      \UnaryInfC{\sFmla{\True}{!A}[m]}
      \DisplayProof
      &
      \AxiomC{\sFmla{\False}{\Diamond !A}[n]}
      \RightLabel{\TRule{\False}{\Diamond}}
      \UnaryInfC{\sFmla{\False}{!A}[m]}
      \DisplayProof\\
      $m$ é novo & $m$ é usado
      \\[1ex]
      \hline}{}
    \end{tabular}
  \end{center}
  \caption{Regras simplificadas para \Log{S5}.}
  \ollabel{tab:rules-S5}
\end{table}

\begin{ex}
  Apresentamos um tableau fechado simplificado que mostra
  $\Log{S5} \Proves \Ax{5}$, O.s.\@, $\Diamond!A \lif \Box\Diamond!A$.
  \begin{oltableau}
    [\pFmla{\False}{\Diamond\formula{A} \lif \Box\Diamond \formula{A}}{1},
      just = \TAss
      [\pFmla{\True}{\Diamond \formula{A}}{1}, just = {\TRule{\False}{\lif}[1]}
        [\pFmla{\False}{\Box\Diamond \formula{A}}{1},
          just = {\TRule{\False}{\lif}[1]}
          [\pFmla{\False}{\Diamond \formula{A}}{2},
            just = {\TRule{\False}{\Box}[3]}
            [\pFmla{\True}{\formula{A}}{3},
              just = {\TRule{\True}{\Diamond}[2]}
                [\pFmla{\False}{\formula{A}}{3},
                  just = {\TRule{\False}{\Diamond}[4]}, close]
            ]
          ]
        ]
      ]
    ]
  \end{oltableau}
\end{ex}

\end{document}
